\chapter*{Conclusion générale}
\addcontentsline{toc}{chapter}{CONCLUSION GÉNÉRALE}

Le projet \textbf{AutoMaze AI} s'inscrit dans une démarche de conception et de développement d'un système de \textbf{classification d'objets routiers par Machine Learning}, destiné à guider les décisions d'un véhicule autonome virtuel. À travers les différentes étapes déjà amorcées, ce travail a permis de poser des bases solides, en définissant clairement les objectifs et le périmètre du projet, en étudiant et en préparant les données, et en concevant une architecture cohérente.

Les choix effectués reposent sur le jeu de données \textbf{COCO 2017}, sur des recadrages carrés de 64~×~64~px et sur un découpage par photo qui évite toute fuite de données. Les trois classifieurs du sujet, \textbf{SVM}, \textbf{Random Forest} et \textbf{k-NN}, seront comparés sur des descripteurs HOG à l'aide du F1 macro. Les prédictions du modèle seront ensuite exploitées par un véhicule virtuel circulant dans les rues de Casablanca, la position des objets étant fournie par le simulateur et leur nature uniquement par le modèle.

Le bilan de faisabilité est encourageant. Les données sont disponibles et préparées, les premiers classifieurs dépassent nettement le hasard, et les moyens nécessaires sont entièrement gratuits. Le projet nous semble donc réalisable dans les délais, sous réserve de l'accord de notre encadrante sur les trois points présentés à la section~\ref{sec:criteres-go}.

Pour la suite, le travail portera sur le réglage des modèles sur la validation, puis sur l'évaluation du modèle final sur le jeu de test et sur de nouvelles images. Si le temps le permet, une attention particulière sera également accordée à l'amélioration des règles de décision, en tenant compte de la distance et de la position des objets, ainsi qu'à l'étude de l'apport d'un CNN entraîné de zéro.

En somme, ce dossier représente une \textbf{phase essentielle de préparation et de structuration} avant la réalisation complète du projet. Il constitue un socle solide pour la mise en œuvre d'un classifieur fiable et pour la démonstration de son utilité dans un contexte de conduite autonome.
